% Guía UVG (pág. 14): traducción al inglés del resumen, con los mismos lineamientos.

Residential projects developed by homebuilders in Guatemala City generate a volume of leads that exceeds the sales team's capacity to attend to them, with no structured criterion for deciding whom to contact first. Under random selection, an average of 31 leads must be contacted to obtain one conversion, a figure that reflects the cost of operating without prioritization.

This work developed a commercial prioritization system built from three components: a model that estimates each lead's probability of conversion, a segmentation that profiles prospective customers, and a Power BI dashboard that delivers both outputs to the sales team. The sales process was first documented using BPMN notation to locate the exact point in the workflow where prioritization needed to be inserted.

The central result is a substantial reduction in commercial effort per conversion. When prioritized by the model, the top 5\% of the portfolio by estimated probability captures 223 of the 739 historical conversions in the analyzed universe, equivalent to 5 leads contacted per conversion obtained, compared to 31 under random selection. At a broader operating point, the model recovers 70\% of actual conversions while contacting less than a third of the portfolio. The segmentation identified six customer profiles, two of which convert at a rate significantly above average and share the same stated purchase motivation.

The finding with the greatest practical implication is that a lead's acquisition channel predicts conversion better than its demographic profile, suggesting that commercial effort should be reoriented toward the contact's origin. The system reorders the portfolio without guaranteeing individual buyers; it is therefore recommended as a criterion for ordering outreach, alongside strengthening the recording of channel variables and validating its impact through a controlled trial before a generalized adoption.
